% Contact sheet for the 23 Oct review figures: every figure on its own page, at print size (the thesis's text
% width), with its one-line message. Build: pdflatex contact_sheet.tex (in this folder, or with \graphicspath).
\documentclass[12pt,a4paper]{article}
\usepackage[margin=1in]{geometry}
\usepackage{graphicx}
\usepackage{amsmath}
\usepackage{xcolor}
\graphicspath{{./}{/home/wesley/multi_drone_control/docs/figures/review_2026-10-23/}}
\pagestyle{plain}
\setlength{\parindent}{0pt}
\newcommand{\fig}[4]{%
  \section*{#1}
  \textit{#2}\par\smallskip
  {\small\color{black!60}#3}\par\vspace{1.2em}
  \includegraphics{#4}%
  \clearpage}
\begin{document}

\fig{1. U: the magnet caught off its slot}
{Radially inward, folding in at the measured catch point keeps the fold-in tilt under $2.7^\circ$ to 16\,cm, while
assuming the slot passes the $5^\circ$ bar from 8\,cm. Along the rim the order reverses past 8\,cm.}
{Preliminary: 6 runs owed (two re-flies, the 20\,cm extension); U re-flies whole on the final tree.}
{fig_u_catch_offset.pdf}

\fig{2. C: our planner against the network on the 0.4\,m/s line}
{Our planner keeps the ring within 3\,mm of the line. The network (our port of Quan et al., best gains) wanders up to
28\,mm and turns further short of the ends.}
{Preliminary: our planner's line has one run; its second and the claim repeats are owed.}
{fig_c_line.pdf}

\fig{3. C: the ring's yaw}
{Our planner holds the ring's yaw within $0.2^\circ$. At its best gains the network lets it swing $22$--$51^\circ$ in
hover and $14^\circ$ on the 0.4\,m/s line, still growing at 60\,s. As ported, the swing grows until the ring
capsizes.}
{Preliminary: our planner's hover is B's baseline hover (same world and plant) until C's own hover cells fly.}
{fig_c_yaw.pdf}

\fig{4. S: solve time against fleet size}
{Fewer SQP iterations per tick let the planner solve bigger fleets within the 60\,ms budget: 4, 5 and 6 quads at
5, 2 and 1 iterations (idle laptop).}
{Preliminary: the bench before the speed-ups; bar B (with the speed-ups and the rig factor) replaces it.}
{fig_s_solve_time.pdf}

\fig{4b. S: plan quality at 8 quads}
{At 8 quads, one iteration per tick leaves the plan outside every quality bar, and two iterations miss the
feed-forward bar by 0.016\,m\,s$^{-2}$. Five pass them all, so the 8-quad claims fly at five.}
{Final for Q6 (the untimed $n=8$ bench).}
{fig_s_quality_n8.pdf}

\fig{5. S: the 8\,$\to$\,4\,$\to$\,8 chain}
{Eight quads carry the 0.9\,kg, 0.5\,m ring through four unannounced releases down to four and four rejoins back to
eight. The ring's tilt peaks at $3.35^\circ$.}
{Preliminary: one exploratory run at 1 iteration; the claim runs at 5 iterations are owed.}
{fig_s_chain.pdf}

\fig{6. The rejoin of a departed carrier}
{A carrier released from the T flies back to its own plate, re-welds 17.8\,s after the rejoin command and folds back
in. The ring's tilt peaks at $1.45^\circ$ in the fold-in (bar $5^\circ$).}
{Final data (E2 passes its frozen bars).}
{fig_rejoin.pdf}

\fig{7a. The pair's swing on the rig, ema against slope}
{With the ema filter, the 0.55\,m pair's 1.4\,Hz see-saw grows after lift-off until the flight is cut. With the slope
filter it stays about half as large and does not grow.}
{Final (rig, 22 Oct). One flight was dropped ($\times$); two were landed by the operator ($\circ$).}
{fig_pair_rig.pdf}

\fig{7b. The pair's swing in the twin}
{In the twin the pair swings only when the throttle is delayed by about 50--60\,ms. Delay on the rates, the rig's
measured timing and the tip hinge stay calm, and both fixes remove the 60\,ms throttle swing.}
{Final (rounds A--D).}
{fig_pair_twin.pdf}

\fig{8. R: what changed at 1\,m}
{At 1\,m the 0.3\,m/s orbit's clock-free path error crosses the 25\,mm envelope. Most of the rise came with the code;
the longer cable adds about 2.6\,mm.}
{Final (the attribution re-fly is in).}
{fig_refly_orbit.pdf}

\fig{9a. Re-render at 1\,m: one flight from the floor}
{From the floor at 1\,m: creep, pretension, lift to 1.0\,m, hold, LAND and unwind. The rods hold their planned
$60^\circ$ through the hold.}
{Re-render of the thesis figure (ring blue instead of gold).}
{fig_flight_1m.pdf}

\fig{9b. Re-render at 1\,m: circle and figure-8}
{At 1\,m, three quads carry the ring around the circle and the figure-8 at 0.125\,m/s about a centimetre from the
reference.}
{Re-render of the thesis figure (ring blue instead of gold; the figure-8 cropped to its path).}
{fig_traj_1m.pdf}

\fig{9c. Re-render at 1\,m: an unannounced detach}
{When one quad of four lets go unannounced at 1\,m, the planner hands its load to the three survivors. The ring's
tilt peaks at $5.9^\circ$ and is back in band in 3.4\,s.}
{Re-render of the thesis figure (ring blue instead of gold).}
{fig_detach_1m.pdf}

\end{document}
